\printedpage{69}
% head: A. Grothendieck : Dérivateurs, Ch. XVIII, Édition des Universités de Montpellier II et Paris VII

\[
\begin{tikzcd}[column sep=small]
 & i(y) \arrow[dr, dotted, "p_\alpha q_y=p_x"] & \\
Z=i(y) \arrow[ur, "\mathrm{id}"] \arrow[dr, "q_y"'] & & X \ , \\
 & X_\alpha \arrow[ur, dotted, "p_\alpha"'] &
\end{tikzcd}
\]
we therefore have, setting $x=(y,p_\alpha q_y)\in\Ob C/X$,
\[
\underbrace{u_x\circ\mathrm{id}_{i(y)}}_{u_x} \ = \ u_\alpha q_y \ ,
\]
so to prove $(**)$, it remains to verify
\[
u_x \ = \ \underbrace{(up_\alpha)q_y}_{u(p_\alpha q_y)\,=\,up_x} \ ,
\]
\msp{page 197}
which is indeed verified by definition of $u_x$. This finishes proving $2^\circ$.

Let us prove $3^\circ$. Since $i'$ is fully faithful, by $1^\circ$, to prove that $i'$ is strictly generating, it suffices to establish that for every object $X$ of $\mathcal{M}$, the family of arrows $i'(x')\longrightarrow X$, indexed by $\Ob C'/X$, is a strict epimorphic family. Now this family contains the family of $\underbrace{i'(\alpha(x))}_{=\,i(x)}\longrightarrow X$, indexed by $\Ob C/X$. Thus it follows from the hypothesis on $i$, and from $2^\circ$, that this family is a strict epimorphic family. We've won! Whew \dots

\textsc{Proof of 4.8.1.} We have
\[
\mathrm{c}' \ \Longrightarrow \ \mathrm{c} \ \Longrightarrow \ \mathrm{b} \ \Longrightarrow \ \mathrm{a} \, ,
\]
the first implication tautological, the second resulting from the fact that $\pi_0$ a regular cardinal for the fcga under consideration, implies that
\[
\mathcal{M} \ \xleftarrow{\ \sim\ } \ \Ind_{\pi_0}(\mathcal{F}^{\pi_0}) \ ,
\]
$\mathcal{F}^{\pi_0}$ equivalent to a small category, whence b).

\msp{page 198}
Finally b) $\Longrightarrow$ a) follows from the fact that $\Ind_{\pi_0}(C)$ is stable under $\varinjlim$ large relative to $\pi_0$ (4.4.15), and that $C$ is a strictly generating subcategory of $\Ind_{\pi_0}(C)$, made up of $\pi_0$-accessible objects. We need to spell out the

\textbf{Lemma 4.8.4.} \textit{Let $C$ be a $\mathfrak{U}$-category, $\pi_0$ an infinite cardinal. Then $C$ is strictly generating in $\mathcal{M}=\Ind C$ and in $\mathcal{M}=\Ind_{\pi_0}C$.}

\msp{page 199}
Indeed, the canonical functor $\mathcal{M}\longrightarrow C^\wedge$ is fully faithful, and we apply criterion c) of 4.8.2.

To finish proving 4.8.1, it therefore remains to establish the implication
\[
\mathrm{a} \ \Longrightarrow \ \mathrm{c}' \, .
\]

\printedpage{70}
% head: A. Grothendieck : Dérivateurs, Ch. XVIII, Édition des Universités de Montpellier II et Paris VII

Let us suppose then that $\mathcal{M}$ satisfies $\mathrm{L}_{\pi_0}$, and admits a small strictly generating subcategory $S$, made up of $\pi_0$-accessible objects, hence a canonical functor
\[
\Ind_{\pi_0}S \ \longrightarrow \ \mathcal{M} \ ,
\]
extending the inclusion $S\hookrightarrow\mathcal{M}$, so as to commute with the $\varinjlim$ large relative to $\pi_0$, and because $S\subset\mathcal{M}_{\pi_0}$, this functor is \textit{fully faithful.} But how to prove that it is essentially surjective? This will not be true, moreover, without a further hypothesis on $S$. The natural idea is to saturate $S$ under the $\varinjlim$
\msp{page 200}
of cardinality $\leqslant\pi_0$ representable in $\mathcal{M}$, and in particular under direct factors, but even then (and even admitting that $\mathcal{M}_{\pi_0}$ itself is equivalent to a small category, and taking $C=\mathcal{M}_{\pi_0}$), it is nevertheless not clear that the fully faithful functor $\Ind_{\pi_0}C\longrightarrow\mathcal{M}$ is essentially surjective, which would also amount to saying that for every $F$ in $\mathcal{M}$, $C/F$ is large relative to $\pi_0$. It is in fact not even clear that $C/F$ is filtering.

So in the end, I doubt that we have a) $\Longrightarrow$ c$'$) or just a) $\Longrightarrow$ b), the situation is visibly not well understood.

But let us prove all the same
\[
\mathrm{b} \ \Longrightarrow \ \mathrm{c}' \, .
\]
For this, we may assume $\mathcal{M}=\Ind_{\pi_0}C$, $C$ small. We then know that the filtration of $\mathcal{M}$ by the $(\mathcal{M}_\pi)_{\pi\geqslant\pi_0}$ is an fcga (prop. 4.7.7 and definition 4.7.8), and even an fcg that is \textit{canonical.} Moreover, I say that the cardinal $\pi_0$ is \textit{useful} and even regular for this filtration (which will finish proving c$'$). We have a slightly more precise result:

\msp{page 201}
\textbf{Corollary 4.8.5.} \textit{Let $C$ be a category equivalent to a small category, $\pi_0$ an infinite cardinal, $\mathcal{M}=\Ind_{\pi_0}(C)\subset C^\wedge$. Then:}
\begin{enumerate}
\item[\textit{1$^\circ$)}] $\pi_0$ \textit{is a cardinal} useful \textit{and even regular for the canonical fcg} $(\mathcal{M}_\pi)_{\pi\geqslant\pi_0}$ \textit{of} $\mathcal{M}$.
\item[\textit{2$^\circ$)}] $\mathcal{M}_{\pi_0}=C^\natural$ \textit{(full subcategory of $\mathcal{M}=\Ind_{\pi_0}C$ made up of the direct factors in $\mathcal{M}$ of objects of $C$), and $\mathcal{M}_{\pi_0}$ is a Karoubi envelope of $C$.}
\end{enumerate}
(\footnote{\textbf{Corollary 4.8.5.1.} \textit{If $C$ is Karoubian, then $\mathcal{M}_{\pi_0}\subset C^\wedge$ is the category made up of the functors} representable \textit{in} $C\overset{\sim}{\hookrightarrow}\mathcal{M}_{\pi_0}$.}, \footnote{\textbf{4.8.5.2.} Spell this out too: Let $u : C\longrightarrow C'$ be a functor, with $C$, $C'$ equivalent to small categories. Equivalent conditions:
\begin{enumerate}
\item[a)] $u$ induces $\Kar C\xrightarrow{\ \sim\ }\Kar C'$.
\item[a$'$)] $u$ fully faithful, and every object of $C'$ is isomorphic to a direct factor of an object of the essential image.
\item[b)] $\Ind(u) : \Ind(C)\longrightarrow\Ind(C')$ is an equivalence.
\item[c)] $\Ind_{\pi_0}(u) : \Ind_{\pi_0}(C)\longrightarrow\Ind_{\pi_0}(C')$ is an equivalence ($\pi_0$ a given infinite cardinal).
\item[d)] $u_! : C^\wedge\longrightarrow C'^\wedge$ is an equivalence.
\item[d$'$)] $u_* : C'^\wedge\longrightarrow C^\wedge$ is an equivalence \texttt{[rather} $u_* : C^\wedge\longrightarrow C'^\wedge$\texttt{]}.
\item[d$''$)] $u^* : C'^\wedge\longrightarrow C^\wedge$ is an equivalence.
\end{enumerate}}).

\printedpage{71}
% head: A. Grothendieck : Dérivateurs, Ch. XVIII, Édition des Universités de Montpellier II et Paris VII

\textsc{Proof.}

$1^\circ$) We already know that
\[
\Ind_{\pi_0}\mathcal{M}_{\pi_0} \ \longrightarrow \ \mathcal{M} \tag*{$(*)$}
\]
is fully faithful, and to see that it is essentially surjective, we note that we have
\[
\mathcal{M} \ \simeq \ \Ind_{\pi_0}C \ \hookrightarrow \ \Ind_{\pi_0}\mathcal{M}_{\pi_0} \ \longrightarrow \ \mathcal{M} \ ,
\]
the composite being isomorphic to the identity functor, is essentially surjective, hence so is $(*)$. This therefore proves that $\pi_0$ is a useful cardinal for the filtration under consideration, let us prove that this cardinal is even regular, hence that for
\[
F \ \in \ \Ob\mathcal{M}_\pi \ , \qquad \text{with } \pi>\pi_0,\ \pi^{\pi_0}\geqslant\pi_1
\]
we have
\[
F \ = \ \varinjlim_I x_i \ , \qquad x_i\in\Ob\mathcal{M}_{\pi_0},\ I \text{ large relative to } \pi_0,\ \card I\leqslant\pi^{\pi_0}.
\]
\msp{page 202}
For this, we use (4.7.6.2), page 171,
\[
\mathcal{M}_\pi \ \subset \ \Filt^{\pi^{\pi_0}} \qquad (\text{if } \pi\geqslant\pi_0,\ \pi^{\pi_0}\geqslant\pi_1),
\]
and the definition of $\Filt^{\pi'}$ (for $\pi'\geqslant\pi_0$) shows us that this is the conclusion sought.

Let us then prove $2^\circ$). Since $\mathcal{M}_{\pi_0}$ is stable under the finite direct $\varinjlim$ which are representable in $\mathcal{M}$ (SGA 4 I 9.9), it is stable under direct factors, and since it contains $C$, we therefore have $C^\natural\subset\mathcal{M}_{\pi_0}$. The reverse inclusion $\mathcal{M}_{\pi_0}\subset C^\natural$ is seen in the standard way, by writing an object $F$ of $\mathcal{M}_{\pi_0}$ as $\varinjlim\limits_I x_i$, the $x_i$ in $C$, $I$ large relative to $\pi_0$, then $F$ is a direct factor of one of the $x_i$.

The fact that $C^\natural$ is a \textit{Karoubi envelope} of $C$ in $\mathcal{M}\subset C^\wedge$ comes from the fact that $\mathcal{M}$ is stable under direct factors, and that $C$ is a full subcategory of $\mathcal{M}$, cf. SGA 4 I (\,?\,\dots).

\textbf{4.8.6. Open questions.} We have thus proved 4.8.1 revised and corrected. There remains the question:
\begin{enumerate}
\item[1)] Give an example where 4.8.1 a) holds, but not b) etc. (\footnote{It suffices to find a nonregular cardinal, for the canonical fcg of an $\mathcal{M}=\Ind_{\pi_0}C$.}).
\item[2)] Bring out serviceable additional conditions which, joined to a), imply b), hence c, c$'$.
\end{enumerate}

\msp{page 203}
Let us also note for the record the following questions.
\begin{enumerate}
\item[3)] Let $\mathcal{M}=\Ind_{\pi_0}(C)$, and $\pi>\pi_0$. Then is $\Filt^\pi$ stable under direct factors? Is this true at least if $\pi=\pi^{\pi_0}$, or what amounts to the same thing (cf. 4.7.3), do we have in this case $\Filt^\pi=\Filt'^{\pi}$? If so, it would be unnecessary in the description 4.7.7 of cardinal filtrations, in the bound on cardinals (4.7.7.1), $\card I\leqslant\pi'^{\pi}$, to put in the restrictive hypothesis $\pi'^{\pi}\geqslant\pi_1$.
\end{enumerate}

To have an affirmative answer to this question, for a given $\pi$, it would suffice (and probably be necessary) to have an affirmative answer to this one: Let $\mathcal{D}$, $\mathcal{D}_0$ be two categories large relative to $\pi_0$, $\mathcal{D}_0$ a direct factor of $\mathcal{D}$, $\mathcal{D}$ admitting a cofinal category $I$ large relative to $\pi_0$, of cardinality $\leqslant\pi$. Then is the same true of $\mathcal{D}_0$? If $p$ is a retraction of $\mathcal{D}$ onto $\mathcal{D}_0$, the difficulty is that I do not know whether $p$ is a cofinal functor. In fact, I see that
\printedpage{72}
% head: A. Grothendieck : Dérivateurs, Ch. XVIII, Édition des Universités de Montpellier II et Paris VII
\msp{page 204}
this is not always the case, even if $\mathcal{D}$ is finite. Let us take e.g. for $\mathcal{D}_0$ a category having final object $e$, and an arrow $x\longrightarrow e$ which is not an isomorphism, and admitting a section (e.g. $\mathcal{D}_0$ the full subcategory of $\boldsymbol{\Delta}$ formed by $\boldsymbol{\Delta}_0$, $\boldsymbol{\Delta}_1$: \begin{tikzcd}[cramped] \boldsymbol{\Delta}_0 \arrow[r, shift left, "\partial_0"] \arrow[r, shift right, "\partial_1"'] & \boldsymbol{\Delta}_1 \arrow[l, bend left=60, "\kappa"] \end{tikzcd} ). Let $\mathcal{D}$ be obtained from $\mathcal{D}_0$ by adding a final object $\tilde{e}$. Then $\mathcal{D}_0$, $\mathcal{D}$ are filtering, and we get a retraction of $\mathcal{D}$ onto $\mathcal{D}_0$, sending $\tilde{e}$ to $x$. This functor is not cofinal, for a cofinal functor must send $\tilde{e}$ to $e$.

More generally, if $\mathcal{D}_0$ is a category such that there exists an arrow
\[
\alpha : e_{\mathcal{D}_0^\wedge} \ \longrightarrow \ x_0 \tag*{$(*)$}
\]
in $\mathcal{D}_0^\wedge$, with $x_0\in\Ob\mathcal{D}_0$, then we see that $\mathcal{D}_0$ is large relative to $\pi_0$ whatever the infinite cardinal $\pi_0$, and that on the other hand, again denoting by $\mathcal{D}$ the category obtained from $\mathcal{D}_0$ by adding a final object to it, there is a retraction of $\mathcal{D}$ onto $\mathcal{D}_0$ which sends $e_{\mathcal{D}}$ to $x_0$. This retraction is cofinal only if $x_0$ is a final object of $\mathcal{D}_0$, i.e. if and only if $(*)$ is an isomorphism. Let us suppose this is not so. But there is a cofinal functor from the subcategory $I_0$ of $\mathcal{D}_0$, such that $\Ob I_0=\{x_0\}$, $\Fl I_0=\End_{I_0}(x_0)=\{p,\mathrm{id}_{x_0}\}$, where $p$ is the composite $x_0\xrightarrow{\text{can.}}e_{\mathcal{D}_0^\wedge}\xrightarrow{\alpha}x_0$. So no counterexample this way!

\msp{page 205}
\begin{enumerate}
\item[4)] Let us consider $\mathcal{M}=\Ind_{\pi_0}(C)$, $C$ small, and let us consider the canonical filtration $(\mathcal{M}_\pi)_{\pi\geqslant\pi_0}$. Give examples where there are non-useful cardinals, or better, where for every cardinal $\pi$, exists a cardinal $\pi'>\pi$ which is non-useful. (This will give the example asked for in $1^\circ$.) On the other hand, give examples of useful cardinals which are not regular -- if possible, as before, such cardinals as large as we want.
\end{enumerate}

\textbf{Definition 4.8.7.} Let $\mathcal{M}$ be a $\mathfrak{U}$-category, $\pi_0$ an infinite cardinal. We say that $\mathcal{M}$ is $\pi_0$-\textit{accessible} if $\mathcal{M}$ is equivalent to a category $\Ind_{\pi_0}(C)$, $C$ a small category, i.e. if it satisfies the equivalent conditions b, c, c$'$ of 4.8.1 (implying condition a), probably strictly weaker). We say that $\mathcal{M}$ is accessible if there exists a cardinal such that $\mathcal{M}$ is $\pi_0$-accessible (\textbf{N.B.} if $\pi_0$ is such a cardinal, then every cardinal $\pi>\pi_0$ such that $\pi^{\pi_0}=\pi$, in particular the cardinals $\pi'^{\pi_0}$, with $\pi'\geqslant 2$, have the same property), or what
\msp{page 206}
amounts to the same thing by 4.8.1, if $\mathcal{M}$ admits an fcga, or again a canonical fcg.

\textbf{Comment 4.8.8.} This definition is in contradiction with the terminology introduced in SGA 4 I 9.4 (p. 114); we shall henceforth designate by the names ``weakly $\pi_0$-accessible category'', ``weakly accessible category'' the notions introduced in 9.4 -- i.e. the categories admitting a full generating subcategory (not even strictly generating) whose objects are $\pi_0$-accessible (resp. accessible).

Let us note that if $\mathcal{M}$ is $\pi_0$-accessible (resp. accessible), the condition of 4.8.1 tells us that $\mathcal{M}$ is weakly $\pi_0$-accessible (resp. weakly accessible), and even that there exists a small full \textit{strictly} generating subcategory, which is made up of $\pi_0$-accessible objects (resp. accessible). I very much doubt that this last necessary condition is also sufficient for
\printedpage{73}
% head: A. Grothendieck : Dérivateurs, Ch. XVIII, Édition des Universités de Montpellier II et Paris VII
$\mathcal{M}$ to be $\pi_0$-accessible (resp. accessible). But by 4.6.3.3 (p. 150), this is so if we assume $C$ stable under small $\varinjlim$ , or what amounts to the same thing, assuming $\mathrm{L}_{\pi_0}$, stable under direct limits of cardinality $\leqslant\pi_0$.

\msp{page 207}
\textbf{Proposition 4.8.9.} \textit{Let $\mathcal{M}$ be an accessible $\mathfrak{U}$-category (4.8.7), $\pi_0$ a cardinal such that $\mathcal{M}$ satisfies $L_{\pi_0}$. Equivalent conditions:}
\begin{enumerate}
\item[\textit{a)}] \textit{$\mathcal{M}$ is a paratopos category (cf. 4.6.3, $3^\circ$, p. 150).}
\item[\textit{b)}] \textit{$\mathcal{M}$ is stable under small direct limits.}
\item[\textit{b$'$)}] \textit{$\mathcal{M}$ is stable under direct limits of cardinality $\leqslant\pi_0$.}
\item[\textit{b$''$)}] \textit{(If $\mathcal{M}$ is equivalent to a categories} $\Ind_{\pi_0}C$\textit{, with $C$ a small category} that is Karoubian.\textit{) The category $C$ is stable (in itself) under direct limits of cardinality $\leqslant\pi_0$.} % [sic] une catégories -> une catégorie, accord du nombre
\item[\textit{c)}] \textit{$\mathcal{M}$ is stable under small inverse limits.}
\item[\textit{c$'$)}] \textit{(Under the hypothesis $\mathcal{M}\simeq\Ind_{\pi_0}(C)$, and when we give ourselves a cardinal $\pi$ such that $\pi=\pi^{\pi_0}\geqslant\pi_1=\sup(\pi_0,\card\textup{réd}\,C)$.) $\mathcal{M}_\pi$ is stable under finite $\varprojlim$, and $\mathcal{M}$ is stable under small infinite products (or again, under filtering $\varprojlim$).}
\item[\textit{c$''$)}] \textit{There exists an infinite cardinal $\pi$, and a small category $\mathcal{D}$ stable under finite $\varprojlim$, such that $\mathcal{M}\simeq\Ind_\pi(\mathcal{D})$, moreover $\mathcal{M}$ stable under infinite products (or again under filtering $\varprojlim$).}
\end{enumerate}
(\footnote{\textbf{N.B.} The equivalence of c$'$, c$''$ with the other conditions is not established (here, but only in 4.11 below.)}, \footnote{\textbf{N.B.} The criteria c$'$, c$''$ seem to me more technical and less useful \texttt{[c$'$, c$''$ above are put in parentheses]}.}).

\textsc{Proof.} The following implications are tautological, or already $\pm$ known:
\[
\begin{tikzcd}[column sep=large, row sep=large]
\mathrm{a} \arrow[r, Rightarrow] & \mathrm{b} \arrow[r, Rightarrow, "\text{cf. 4.4.14}"] \arrow[dl, Rightarrow] & \mathrm{b}' \arrow[r, Rightarrow, "\bullet"] & \mathrm{b}'' \\
\mathrm{c}\ ( \arrow[r, Rightarrow, "\bullet"] & \mathrm{c}' \arrow[r, Rightarrow] & \mathrm{c}'')\ , &
\end{tikzcd}
\]
\msp{page 208}
where only the dot-marked implications b$'\Longrightarrow$ b$''$ and c $\Longrightarrow$ c$'$ call for a proof. For b$'\Longrightarrow$ b$''$, we note that the hypothesis $C$ Karoubian implies that $C\simeq\mathcal{M}_{\pi_0}$ (4.8.5, $2^\circ$), now we know that $\mathcal{M}_{\pi_0}$ is stable in $\mathcal{M}$ under direct limits of cardinality $\leqslant\pi_0$ (and these are representable in $\mathcal{M}$ by b$'$), hence $\mathcal{M}_{\pi_0}$ admits $\varinjlim$ of the type under consideration, hence $C$ too.

c $\Longrightarrow$ c$'$. Let us note that if $\mathcal{M}=\Ind_{\pi_0}(C)$, then the $\varprojlim$ representable in $\mathcal{M}$ are computed in $C^\wedge$. It follows that if $\pi\geqslant\pi_1$, then $\Filt^\pi$ is stable under finite $\varprojlim$ ($\Filt^\pi$ then being given by condition a$''$ of 4.7.1), on the other hand, if in addition $\pi^{\pi_0}=\pi$, $\pi$ is a useful cardinal, hence $\mathcal{M}\simeq\Ind_\pi(\mathcal{D})$, where $\mathcal{D}=\Filt^\pi=\mathcal{M}_\pi$. It now suffices to see that $\Ind_\pi(\mathcal{D})$ is then stable under finite $\varprojlim$, to conclude that the same holds for $\mathcal{M}$, hence that $\mathcal{M}$ is stable under arbitrary inverse limits, if it is under small products or under small filtering $\varprojlim$, whence c$'$. It therefore remains to prove the following lemma. % [sic] petites produites -> petits produits, accord